%%%%%%%%%%  scan idx 309  |  folio 296  %%%%%%%%%%
(Bien entendu, pour \(X_i=Y_i=S\), on retrouve la structure de \(K^{\bullet}(S)\)-module de \(K_{\bullet}(S)\).)

Si \(X\) est plat sur \(S\), on en déduit, pour \(n\geq1\), des homomorphismes
\[
 K^{\bullet}(X/S)^{\otimes n}\longrightarrow K^{\bullet}\bigl((X/S)^n\bigr),
\]
qui devraient en un sens remplacer la structure d'anneau absente sur \(K^{\bullet}(X/S)\).

\subsection*{BIBLIOGRAPHIE}
\begin{enumerate}[label={[\arabic*]}]
\item Deligne, P., papiers secrets.
\item Grauert, H., Ein Theorem der Analytischen Garbentheorie \dots, Pub. Math. IHES, n\textsuperscript{o}5 (1960).
\item Grothendieck, A., Théorèmes de finitude pour la cohomologie des faisceaux, Bull. Soc. Math. France, 84, 1956, p.1-7.
\item[{[V]}] Verdier, J.-L., Catégories dérivées, IHES (1964), miméographié.
\end{enumerate}



% Expose V macros
\providecommand{\Chat}{\widehat C}
\providecommand{\Ghat}{\widehat G}
\providecommand{\Gg}{\widehat G_g}
\providecommand{\Ga}{\widehat G_a}
\providecommand{\Grhat}{\widehat{\operatorname{Gr}}}
\providecommand{\Abfil}{\underline{\mathrm{Ab}}\,\underline{\mathrm{fil}}}
\providecommand{\Kalgfil}{K\text{-}\mathrm{Alg.fil.}}
\providecommand{\Hom}{\operatorname{Hom}}
\providecommand{\Sym}{\mathfrak S}
\renewcommand{\id}{\operatorname{Id}}
\providecommand{\Log}{\operatorname{Log}}
\providecommand{\Toddc}{\operatorname{Todd}}
\renewcommand{\Ob}{\operatorname{Ob}}
\providecommand{\cG}{\mathcal G}
\providecommand{\cCzero}{\mathcal C_0}
\providecommand{\Fil}{\operatorname{Fil}}


